%! AP Statistics — Unit 3, Lesson 3.6 — Activity KEY
\documentclass[10pt]{article}
\usepackage{apstats-article}
\usepackage{apstats-key}

\begin{document}

\pageheader{Unit 3, Lesson 3.6}{Group Activity KEY: Choose the Best Design}
\namepartnerperiod

\vspace{0.1in}

\textbf{Scenario 1:} Math curriculum, grade-level split.
\begin{practicebox}
\small
\begin{enumerate}[leftmargin=1.5em, itemsep=5pt]
  \item \ans{Block design}
  \item \ans{Math level (grade-level vs.\ below grade-level)}
  \item \ans{Math proficiency level is likely to affect test scores. Blocking removes this source of variability, giving a more precise estimate of the curriculum's effect. With 30 in each block, random assignment within blocks is straightforward.}
  \item Diagram: 60 students $\to$ Block 1 (grade-level, 30) / Block 2 (below grade-level, 30); within each block, random assign 15 to new curriculum, 15 to traditional.
\end{enumerate}
\end{practicebox}

\vspace{0.08in}

\textbf{Scenario 2:} Acne treatment, same patient both cheeks.
\begin{practicebox}
\small
\begin{enumerate}[leftmargin=1.5em, itemsep=5pt]
  \item \ans{Matched pairs design}
  \item \ans{Each patient serves as their own block; both cheeks are the same person's skin, controlling for all individual factors (genetics, hygiene, etc.).}
  \item \ans{Randomly assign (e.g., coin flip) which cheek gets Gel A and which gets Gel B.}
  \item Diagram: 40 patients $\to$ each person: left cheek / right cheek; random assignment (coin flip) $\to$ Gel A or Gel B.
\end{enumerate}
\end{practicebox}

\vspace{0.08in}

\textbf{Scenario 3:} Three corn seeds, 4 fields.
\begin{practicebox}
\small
\begin{enumerate}[leftmargin=1.5em, itemsep=5pt]
  \item \ans{Block design}
  \item \ans{Field (soil quality)}
  \item \ans{4 plots per treatment per block (48 plots / 4 fields / 3 treatments = 4)}
  \item Diagram: 4 blocks (fields), each with 12 plots; within each block randomly assign 4 plots to Seed A, 4 to Seed B, 4 to Seed C.
  \item \ans{Removes the field-to-field variability in soil quality from the error, so differences between seeds are estimated more precisely.}
\end{enumerate}
\end{practicebox}

\vspace{0.08in}

\begin{practicebox}
\small
\textbf{Decision Chart:}

\renewcommand{\arraystretch}{1.6}
\begin{tabularx}{\linewidth}{@{} X c @{}}
\textbf{Situation} & \textbf{Best design} \\
\hline
Homogeneous units; no nuisance variable. & \ans{CRD} \\
Nuisance variable known; cannot pair. & \ans{Block design} \\
Units can be paired 1-to-1. & \ans{Matched pairs} \\
Same unit can receive both treatments. & \ans{Matched pairs} \\
\end{tabularx}
\end{practicebox}

\end{document}
